%% volumes.tex — Foundations of Scientific Computing의 권 구성 (모든 권이 함께 쓴다)
%%
%% 부 하나가 권 하나이다. 장 번호는 책 전체에서 이어지므로 "49장"은 어느 권에서든 같은 장이다.
%% 각 권의 본 파일(volN.tex)은 \documentclass 다음에 이 파일을 \input하고 \sdsetvolume{N}을 부른다.
%% 다른 권의 장을 \cref로 가리킬 수 있도록 그 장의 참조를 xr-hyper로 읽는다(링크는 그 권의 PDF로 간다).
%% 읽는 것은 build/xref/의 사본이다: tools/build.sh가 장의 aux에서 biblatex의 refsection 라벨
%% (권마다 1부터 다시 세어 겹친다)을 뺀 \newlabel 줄만 옮겨 둔다.
%% 빌드: tools/build.sh (모든 권을 두 번씩 돌려 권 사이의 참조를 맞춘다)

\makeatletter
% 권 번호, 제목, 장 범위
\newcommand\sd@voltitle[1]{\ifcase#1\or
  Mathematical Preliminaries\or
  Computer Arithmetic and Systems\or
  Numerical Linear Algebra\or
  Approximation and Analysis\or
  Optimization\or
  Probability and Statistical Inference\or
  Statistical Models\or
  Geometry, Graphs, and Graphics\fi}
\newcommand\sd@volrange[1]{\ifcase#1\or
  Chapters 1--7\or Chapters 8--18\or Chapters 19--28\or Chapters 29--41\or
  Chapters 42--48\or Chapters 49--60\or Chapters 61--69\or Chapters 70--78\fi}
\newcommand\sd@volfirst[1]{\ifcase#1\or 1\or 8\or 19\or 29\or 42\or 49\or 61\or 70\fi}
\def\sd@nvolumes{8}

% 장 파일과 그 장이 속한 권. 다른 권의 장은 그 aux를 읽어 참조만 가져온다.
\newcommand\sd@chapters{%
  \sd@ch{1}{01-order}\sd@ch{1}{02-integers}\sd@ch{1}{03-asymptotics}\sd@ch{1}{04-matrix-analysis}%
  \sd@ch{1}{05-complex}\sd@ch{1}{06-recurrences}\sd@ch{1}{07-combinatorics}%
  \sd@ch{2}{08-floating-point}\sd@ch{2}{09-rounding-error}\sd@ch{2}{10-exact-sums}%
  \sd@ch{2}{11-elementary-functions}\sd@ch{2}{12-number-text}\sd@ch{2}{13-architecture}%
  \sd@ch{2}{14-parallel}\sd@ch{2}{15-sorting}\sd@ch{2}{16-data-structures}%
  \sd@ch{2}{17-memory-model}\sd@ch{2}{18-language-boundary}%
  \sd@ch{3}{19-lu}\sd@ch{3}{20-symmetric}\sd@ch{3}{21-qr}\sd@ch{3}{22-least-squares}%
  \sd@ch{3}{23-eigenvalues}\sd@ch{3}{24-svd}\sd@ch{3}{25-matrix-functions}\sd@ch{3}{26-sparse}%
  \sd@ch{3}{27-krylov}\sd@ch{3}{28-large-eigen}%
  \sd@ch{4}{29-polynomial-interpolation}\sd@ch{4}{30-rational-approximation}\sd@ch{4}{31-splines}%
  \sd@ch{4}{32-multivariate-interpolation}\sd@ch{4}{33-numerical-differentiation}\sd@ch{4}{34-roots}%
  \sd@ch{4}{35-quadrature-1}\sd@ch{4}{36-quadrature-2}\sd@ch{4}{37-fourier}\sd@ch{4}{38-fft}%
  \sd@ch{4}{39-special-1}\sd@ch{4}{40-special-2}\sd@ch{4}{41-special-3}%
  \sd@ch{5}{42-optimization-basics}\sd@ch{5}{43-scalar-minimization}%
  \sd@ch{5}{44-line-search-quasi-newton}\sd@ch{5}{45-trust-region}\sd@ch{5}{46-derivative-free}%
  \sd@ch{5}{47-nonlinear-least-squares}\sd@ch{5}{48-constraints-duality}%
  \sd@ch{6}{49-probability-spaces}\sd@ch{6}{50-distribution-families}%
  \sd@ch{6}{51-moments-limit-theorems}\sd@ch{6}{52-multivariate-distributions}%
  \sd@ch{6}{53-random-numbers}\sd@ch{6}{54-estimation}\sd@ch{6}{55-tests-intervals}%
  \sd@ch{6}{56-parametric-tests}\sd@ch{6}{57-nonparametric}\sd@ch{6}{58-contingency-proportions}%
}

% \sdsetvolume{N}: 이 권의 번호를 정하고 다른 권의 장 참조를 읽는다(전문에서 부른다)
\newcommand\sdsetvolume[1]{%
  \gdef\sd@thisvol{#1}%
  \title{Foundations of Scientific Computing}%
  \subtitle{Theory, Algorithms, and Implementation}%
  \gdef\sd@volumeline{Volume \@Roman{#1}\quad\sd@voltitle{#1}}%
  \hypersetup{pdftitle={Foundations of Scientific Computing, Volume \@Roman{#1}: \sd@voltitle{#1}}}%
  \def\sd@ch##1##2{\ifnum##1=#1\relax\else
    \IfFileExists{build/xref/##2.aux}{\externaldocument{build/xref/##2}[vol##1.pdf]}{}\fi}%
  \sd@chapters}

% \sdvolumechapters: 이 권의 장을 차례로 넣는다(본문에서 부른다)
\newcommand\sdvolumechapters{%
  \setcounter{chapter}{\numexpr\sd@volfirst{\sd@thisvol}-1\relax}%
  \def\sd@ch##1##2{\ifnum##1=\sd@thisvol\relax\expandafter\@firstofone\else\expandafter\@gobble\fi{\include{chapters/##2}}}%
  \sd@chapters}

% \sdvolumefront: 속표지, 권 목록, 차례
\newcommand\sdvolumefront{%
  \maketitle
  \frontmatter
  \sdvolumelist
  \tableofcontents}

% 권 목록: 속표지 다음 쪽에 이 책 전체의 권 구성을 적는다
\newcommand\sdvolumelist{%
  \thispagestyle{empty}%
  \null\vskip 2cm
  {\noindent\Large Foundations of Scientific Computing\par}%
  \vskip 0.4em
  {\noindent Theory, Algorithms, and Implementation\par}%
  \vskip 1.5cm
  \begingroup
  \@tempcnta=1
  \loop
    \noindent\hbox to 6em{\ifnum\@tempcnta=\sd@thisvol\bfseries\fi Volume \@Roman\@tempcnta\hfil}%
    \hbox to 17em{\ifnum\@tempcnta=\sd@thisvol\bfseries\fi\sd@voltitle{\the\@tempcnta}\hfil}%
    {\ifnum\@tempcnta=\sd@thisvol\bfseries\fi\sd@volrange{\the\@tempcnta}}\par
    \vskip 0.3em
    \advance\@tempcnta by 1
  \ifnum\@tempcnta<\numexpr\sd@nvolumes+1\relax\repeat
  \endgroup
  \vfill
  {\noindent\small Chapters are numbered continuously across all volumes.\par}%
  \newpage}
\makeatother
